\documentclass[11pt,a4paper]{article}
\usepackage[T1]{fontenc}
\usepackage[utf8]{inputenc}
\usepackage{lmodern}
\usepackage[margin=25mm,headheight=14pt]{geometry}
\usepackage{amsmath,amssymb,amsthm,mathtools,bm}
\usepackage{microtype,booktabs,array,longtable}
\usepackage{enumitem}
\usepackage{siunitx}
\usepackage{xcolor}
\usepackage{fancyhdr}
\usepackage[hyphens]{url}
\usepackage{hyperref}
\hypersetup{colorlinks=true,linkcolor=blue!45!black,citecolor=blue!45!black,
 urlcolor=blue!45!black,pdftitle={Critical Transseries at a Moving Fold},
 pdfauthor={Research continuation for the ProveIt project}}
\pagestyle{fancy}
\fancyhf{}
\fancyhead[L]{\small Critical transseries at a moving fold}
\fancyhead[R]{\small Research manuscript}
\fancyfoot[C]{\thepage}
\setlength{\parskip}{4pt}
\setlength{\parindent}{1.2em}
\setcounter{tocdepth}{2}
\setlist[enumerate]{itemsep=4pt,topsep=5pt}
\allowdisplaybreaks[2]
\newtheorem{theorem}{Theorem}[section]
\newtheorem{lemma}[theorem]{Lemma}
\newtheorem{proposition}[theorem]{Proposition}
\newtheorem{corollary}[theorem]{Corollary}
\theoremstyle{definition}
\newtheorem{definition}[theorem]{Definition}
\newtheorem{example}[theorem]{Example}
\newtheorem{question}{Research question}[section]
\theoremstyle{remark}
\newtheorem{remark}[theorem]{Remark}
\newcommand{\C}{\mathbb C}
\newcommand{\N}{\mathbb N}
\newcommand{\Q}{\mathbb Q}
\newcommand{\R}{\mathbb R}
\newcommand{\OO}{\mathcal O}
\newcommand{\supp}{\operatorname{supp}}
\newcommand{\dd}{\,\mathrm d}
\newcommand{\tc}{t_{\mathrm c}}
\newcommand{\zc}{z_{\mathrm c}}
\newcommand{\vc}{v_{\mathrm c}}
\newcommand{\coeff}[2]{[#1^{#2}]}
\newcommand{\repoLabel}[1]{\texttt{\small\detokenize{#1}}}

\makeatletter
\newcommand{\tableinput}[1]{\@@input #1 }
\makeatother

\begin{document}
\begin{titlepage}
\vspace*{13mm}
\begin{center}
{\small\scshape Transseries, analytic inversion, and critical asymptotics}\par
\vspace{10mm}
{\LARGE\bfseries Critical Transseries\par at a Moving Fold}\par
\vspace{7mm}
{\large A uniform two-sheet atlas,\par large-sector crossover, and resonant action splitting}\par
\vspace{12mm}
{\normalsize Research continuation for the ProveIt project}\par
\vspace{3mm}
{\normalsize 29 September 2026}
\end{center}
\vspace{9mm}
\begin{abstract}
The exponential sectors in the inverse of a logarithmic--exponential map
have poles of order $2n-1$ at a critical core value. Their deepest-pole sum
is a square root, but that formal sum does not locate the moving branch
point and is not a uniform approximation at the transition. We resolve
this explicitly posed project-local problem for
\[
 \delta+\log(1+\delta/w)+wz e^{-\delta}=0.
\]
Writing $s=w+1$ and $z=s^2\zeta$, we construct an exact holomorphic two-sheet
atlas on an explicit complex domain. Its discriminant gives
$\zc(s)=s^2/2+s^3/3+5s^4/8+\cdots$ and an all-order polynomial profile
expansion with a geometric error bound valid at the branch point itself.
A uniform coefficient theorem then proves the double-scaling law
$a_n(\tau/n)/a_n(0)\to e^{-2\tau/3}$ and its first correction, where
$a_n(s)=s^{2n-1}b_n(-1+s)$ normalizes the $n$th exponential sector.
Finally, for general analytic folds driven by finitely many exponential
actions, we prove an absolutely convergent, resonance-aware transseries
representation and give a constructive finite positive generating set for
its support. Explicit cancellations can change the first fractional
action or remove the splitting altogether. The arguments use classical
local analytic and coefficient-extraction methods; the contribution is
the quantitative specialization, crossover calculation, and support
construction, not a new general theory of analytic folds.
\end{abstract}
\vfill
\noindent\fbox{\begin{minipage}{0.93\textwidth}
\small\textbf{Scope and verification.}
This manuscript answers Research question 11.8 of the predecessor
\emph{Support-Controlled Reversion and the Exact One-Exponential
Substitution Group} in its independent-parameter local setting.
The new results are not Lean-formalized. Exact symbolic checks and
100-digit numerical diagnostics are supplied, but are not substitutes
for the proofs or for independent review. Worldwide priority is not claimed.
\end{minipage}}
\end{titlepage}

{\small
\setlength{\parskip}{0pt}
\tableofcontents
}
\newpage

\section{Research target, contribution, and boundaries}
\label{sec:intro}

\subsection{The gap between a formal pole sum and a uniform inverse}
The ProveIt transseries program separates several logically different
levels: finite algebraic identities, formal support-controlled operations,
analytic realization, and quantitative error bounds. This separation is
particularly important near a zero of the derivative of the function
being inverted. A coefficientwise expansion can be perfectly correct
while its truncations fail precisely where infinitely many of its
singular coefficients compete.

The inspected repository module \texttt{QuadraticCoreCatalan.lean}
formalizes the finite Catalan algebra underlying the quadratic equation
\[
 A\varepsilon+c\varepsilon^2+\sigma q=0.
\]
The package README and the module itself distinguish that algebra from
the assembled power-series existence statement and from the identification
of deepest poles in special-function inversions
\cite{repo-readme,repo-catalan}. We maintain that distinction here.

The predecessor manuscript \cite{predecessor} applies exact-core inversion
to the logarithmic--exponential example. It proves an exact rational
formula for the sector functions $b_n(w)$ and identifies their leading
poles at $w=-1$. Its Research question 11.8 asks whether the resulting
formal square-root expression can be promoted to a uniform complex
transition approximation, with error bounds and the nearby branch points
actually located. This is the central target of the present article.

The answer is affirmative, but the branch point must move. Keeping the
quadratic model's branch point fixed gives an approximation whose error
is of order $s^{3/2}$ on a natural test path, rather than the $s^2$ order
one might infer from a regular correction expansion. The corrected
coordinate is
\begin{equation}\label{eq:intro-q}
 q^2=1-\frac{z}{\zc(s)},\qquad s=w+1,
\end{equation}
not $1-2z/s^2$.

\subsection{What is established here}
The main results have three complementary forms.

First, an explicit two-sheet analytic representation is proved for
$|s|\le 10^{-2}$ and $|z/s^2|<3/4$. It yields a convergent expansion
\[
 \delta=s\left(\sum_{j\ge0}s^jV_j(q)-1\right),
 \qquad V_j\in\Q[q],\qquad \deg V_j\le j+1,
\]
with a bound uniform in the truncation order and valid for $|q|\le1/2$,
including $q=0$. The first five profiles are given explicitly.

Second, the same atlas controls the limit of large sector index. If
$a_n(s)=s^{2n-1}b_n(-1+s)$, then uniformly for bounded complex $\tau$,
\begin{equation}\label{eq:intro-crossover}
 \frac{a_n(\tau/n)}{a_n(0)}
 =e^{-2\tau/3}
 \left[1+\frac{\tau/3-37\tau^2/36}{n}+\OO(n^{-2})\right].
\end{equation}
Thus the first drift of the branch point has an order-one effect on high
sectors in the window $s\asymp n^{-1}$.

Third, an analytic fold perturbed by finitely many exponentials admits an
absolutely convergent transseries whose leading splitting action is half
the first \emph{noncancelled} action in its discriminant. Equal actions
must be combined before selecting this first term. We give a finite,
explicitly enumerable positive generating set for all subsequent action
gaps. This answers a substantial analytic subclass of the predecessor's
multi-action questions, not its unrestricted formal or resurgent versions.

\subsection{Classical input and the meaning of novelty}
Square-root branching at a simple critical point is classical. The
Lambert $W$ branch expansion is an especially relevant example
\cite{corless}. Extraction of coefficient asymptotics from such a
singularity is also classical \cite{flajolet-odlyzko}. General transseries
background, including the need for appropriate support conditions, can
be found in \cite{edgar}.

We do not claim to replace these theories. Instead, we prove the precise
uniform analytic statement that the earlier formal calculation left open,
compute its moving discriminant and polynomial profiles, and derive a
specific joint critical/large-sector limit. The proofs below are
self-contained at the level of elementary complex analysis, including
Cauchy's formula, Rouch\'e's theorem, the holomorphic implicit-function
theorem, and the gamma-function coefficient formula. A constructive
finite-grid refinement is then added for analytic finite-action folds.

\subsection{An essential parameter qualification}
\label{sec:qualification}
The variable $z$ is an \emph{independent analytic perturbation parameter}.
For the original fixed map $y+\log y+e^{-y}$, putting
$x=w+\log w$ gives the specialization $z=e^{-x}=e^{-w}/w$.
As $w\to-1$ on a local logarithm branch, that value tends to $-e$, not to
zero. Therefore our small-$(s,z)$ critical chart is not, by itself, an
analysis of the fixed-coupling map at $w=-1$.

More generally the equation studied here is the local form of
\[
 y+\log y+\lambda e^{-y}=w+\log w,
 \qquad \lambda=we^wz.
\]
It is a small-coupling, independent-parameter analytic family. Its sector
functions are exactly those of the predecessor, so its singularity
analysis directly controls their critical behavior. Global continuation
to the physical specialization $\lambda=1$ is a different question and
is not claimed in this article.

\section{From the sector equation to a regular critical chart}
\label{sec:chart}

\subsection{The known sector functions}
For $w\ne0,-1$, let $\delta(w,z)$ denote the solution near $z=0$ with
$\delta(w,0)=0$ of
\begin{equation}\label{eq:delta}
 \delta+\log\left(1+\frac{\delta}{w}\right)+wz e^{-\delta}=0.
\end{equation}
The logarithm in this equation is the local branch equal to zero when
$\delta=0$. Write
\begin{equation}\label{eq:sectors}
 \delta(w,z)=\sum_{n\ge1} b_n(w)z^n.
\end{equation}
The predecessor gives
\begin{equation}\label{eq:known-bn}
 b_n(w)=\frac{(-1)^n}{n!}(T-n)^{n-1}\frac{w^{n+1}}{w+1},
 \qquad T=\frac{w}{w+1}\frac{\mathrm d}{\mathrm dw}.
\end{equation}
For reference,
\begin{align}
 b_1(w)&=-\frac{w^2}{w+1},\nonumber\\
 b_2(w)&=-\frac{w^3(2w^2+2w-1)}{2(w+1)^3},\label{eq:first-bn}\\
 b_3(w)&=-\frac{w^4(9w^4+18w^3-11w+1)}{6(w+1)^5}.\nonumber
\end{align}
Its deepest-pole identity is
\begin{equation}\label{eq:old-poles}
 \lim_{s\to0}s^{2n-1}b_n(-1+s)
 =-\frac{(2n-3)!!}{n!},\qquad (-1)!!=1.
\end{equation}
Consequently the most singular part of every sector has the formal sum
\begin{equation}\label{eq:old-sqrt}
 -\sum_{n\ge1}\frac{(2n-3)!!}{n!}\frac{z^n}{s^{2n-1}}
 =s\left(\sqrt{1-\frac{2z}{s^2}}-1\right).
\end{equation}
Equations \eqref{eq:known-bn}--\eqref{eq:old-sqrt} are predecessor results,
not new assertions of priority here.

\subsection{Blowing up the coalescing point}
Set
\begin{equation}\label{eq:coordinates}
 s=w+1,\qquad t=y+1=s+\delta,
 \qquad h(t)=t+\log(1-t)=-\sum_{j\ge2}\frac{t^j}{j}.
\end{equation}
For $|s|,|t|<1$, direct substitution in \eqref{eq:delta} gives
\begin{equation}\label{eq:Z}
 z=Z(s,t):=\frac{e^{t-s}}{1-s}\bigl(h(t)-h(s)\bigr).
\end{equation}
Indeed $1+\delta/w=(1-t)/(1-s)$, so the logarithmic difference in
\eqref{eq:delta} is exactly the local difference used in \eqref{eq:Z}.

The variables $t$ and $s$ vanish at the same scale. Put $t=sv$ and
$z=s^2\zeta$. Define
\begin{equation}\label{eq:Phi}
 \Phi(s,v)=\frac{e^{s(v-1)}}{1-s}
 \frac{h(sv)-h(s)}{s^2}.
\end{equation}
The apparent singularity at $s=0$ is removable because
\begin{equation}\label{eq:H-series}
 \frac{h(sv)-h(s)}{s^2}
 =-\sum_{j\ge2}\frac{v^j-1}{j}s^{j-2}.
\end{equation}
Thus
\begin{equation}\label{eq:Phi-zero}
 \Phi(0,v)=P(v):=\frac{1-v^2}{2},\qquad
 \Phi(s,1)=0,\qquad \Phi_v(s,1)=-\frac1{(1-s)^2}.
\end{equation}
There is a unique inverse germ $v(s,\zeta)$ near $\zeta=0$ satisfying
$v(s,0)=1$. Write
\begin{equation}\label{eq:an}
 v(s,\zeta)=1+\sum_{n\ge1}a_n(s)\zeta^n.
\end{equation}
The analytic dependence on $s$ removes the original coefficient poles:
\begin{equation}\label{eq:normalization}
 b_n(-1+s)=s^{1-2n}a_n(s)\quad(s\ne0),\qquad
 a_n(0)=-\frac{(2n-3)!!}{n!}.
\end{equation}
The latter identity follows from $v(0,\zeta)=\sqrt{1-2\zeta}$, with the
square root equal to one at zero.

\section{An explicit holomorphic two-sheet atlas}
\label{sec:atlas}

\subsection{Quantitative perturbation estimates}
\begin{lemma}\label{lem:bounds}
For $|s|\le r:=1/100$ and $|v|\le2$,
\begin{equation}\label{eq:Phi-bounds}
 |\Phi(s,v)-P(v)|<\frac7{50},\qquad
 |\Phi_v(s,v)+v|<\frac4{25}.
\end{equation}
\end{lemma}
\begin{proof}
Let $H=(h(sv)-h(s))/s^2$ and $E=e^{s(v-1)}/(1-s)$, so $\Phi=EH$.
The series \eqref{eq:H-series} gives
\[
 |H-P|\le R_h:=\frac r3\left(\frac8{1-2r}+\frac1{1-r}\right).
\]
Since $|s(v-1)|\le3r$, the elementary bound
$e^u\le(1-u)^{-1}$ for $0\le u<1$ yields
\[
 |E|\le E_0:=\frac1{(1-3r)(1-r)},\qquad |E-1|\le E_0-1.
\]
Also $|P|\le5/2$, hence
\[
 |\Phi-P|\le\frac52(E_0-1)+E_0R_h
 =\frac{37787705}{279504918}<\frac7{50}.
\]
Differentiating the convergent series for $H$ gives
$|H_v+v|\le4r/(1-2r)$. Since $E_v=sE$,
\begin{align*}
 |\Phi_v+v|
 &\le2(E_0-1)+E_0\left[\frac{4r}{1-2r}
                  +r\left(\frac52+R_h\right)\right]\\
 &=\frac{21177832}{139752459}<\frac4{25}.
\end{align*}
All estimates are absolute complex estimates; no real positivity of $s$
is assumed.
\end{proof}

\begin{lemma}\label{lem:two-roots}
For $|s|\le1/100$ and $|\zeta|<3/4$, the equation
$\Phi(s,v)=\zeta$ has exactly two roots in $|v|<2$, counted with
multiplicity. There is exactly one critical point $\vc(s)$ in $|v|<2$,
and it is simple as a zero of $\Phi_v$. Its critical value is denoted by
$\rho(s)=\Phi(s,\vc(s))$.
\end{lemma}
\begin{proof}
On $|v|=2$,
\[
 |P(v)-\zeta|=\frac12|v^2-(1-2\zeta)|\ge\frac34
 \quad\text{when }|\zeta|\le\frac34.
\]
The first estimate in Lemma~\ref{lem:bounds} and Rouch\'e's theorem compare
$\Phi-\zeta$ with $P-\zeta$. Both roots of the latter lie strictly within
$|v|=2$. The second estimate compares $\Phi_v$ with $-v$ on the same
circle, proving that the total number of critical points, with
multiplicity, is one. The unique zero must consequently be simple.
Holomorphic dependence of $\vc$ follows from the implicit-function theorem.
\end{proof}

\subsection{The critical value and a useful enclosure}
Differentiate \eqref{eq:Z}. The critical-point equation is
\begin{equation}\label{eq:k}
 k(\tc(s))=h(s),\qquad
 k(t):=h(t)+h'(t)=t+\log(1-t)-\frac{t}{1-t}.
\end{equation}
Here $k(0)=0$ and $k'(0)=-1$, so $\tc$ is holomorphic near zero. Moreover,
\begin{equation}\label{eq:rho-exact}
 \vc(s)=\frac{\tc(s)}s,\qquad
 \zc(s)=s^2\rho(s)=\frac{e^{\tc(s)-s}}{1-s}
                          \frac{\tc(s)}{1-\tc(s)}.
\end{equation}
The quotients have removable singularities at $s=0$.

\begin{lemma}\label{lem:rho-enclosure}
For $0<|s|\le1/100$, the critical point satisfies
$|\tc(s)|<|s|^2$. Its normalized critical value obeys
\begin{equation}\label{eq:rho-enclosure}
 \left|\rho(s)-\frac12\right|<\frac1{250}.
\end{equation}
In particular, $\rho$ is nonzero throughout this disk and $|\rho|<0.504$.
\end{lemma}
\begin{proof}
The exact series
\[
 k(t)=-t-\sum_{j\ge2}\frac{j+1}{j}t^j
\]
implies, on $|t|=|s|^2$,
\[
 |k(t)+t|\le\frac{3|s|^4}{2(1-|s|^2)},\qquad
 |h(s)|\le\frac{|s|^2}{2(1-|s|)}.
\]
The sum of these bounds is strictly smaller than $|s|^2$ for
$|s|\le1/100$. Rouch\'e's theorem applied to $k(t)-h(s)$ and $-t$
therefore proves the asserted location. This solution is the critical
point of Lemma~\ref{lem:two-roots} because $|\tc/s|<|s|<2$.

Dividing \eqref{eq:k} by $s^2$ and separating the first term of $h(s)$
also gives
\begin{equation}\label{eq:tc-half-bound}
 \left|\frac{\tc(s)}{s^2}-\frac12\right|
 \le\frac{|s|}{3(1-r)}+\frac{3|s|^2}{2(1-r^2)}.
\end{equation}
In the exponential factor of \eqref{eq:rho-exact}, write
\[
 \frac{e^{\tc-s}}{1-s}=\exp(\tc-s-\log(1-s)).
\]
The exponent has absolute value at most
$L=r^2(1+1/[2(1-r)])$. Put $E_c=(1-L)^{-1}$. The same exponential
estimate gives a bound $E_c$ for this factor and $E_c-1$ for its
difference from one. Substitution of \eqref{eq:tc-half-bound} into
\eqref{eq:rho-exact} now bounds $|\rho-1/2|$ by
\[
 \left[\frac r{3(1-r)}+\frac{3r^2}{2(1-r^2)}\right]
       \frac{E_c}{1-r^2}
 +\frac12\left(\frac{E_c}{1-r^2}-1\right)
 =\frac{21851394853}{5997897210294}<\frac1{250}.
\]
The estimate extends to $s=0$ by continuity.
\end{proof}

\subsection{Exact preparation by contour moments}
\begin{theorem}[Uniform two-sheet atlas]\label{thm:atlas}
There are holomorphic functions $A(s,\zeta)$ and $B(s,\zeta)$ on
\[
 |s|<\frac1{100},\qquad |\zeta|<\frac34,
\]
with $B$ nowhere zero, such that the two roots in $|v|<2$ are exactly
\begin{equation}\label{eq:atlas}
 v_\pm(s,\zeta)=A(s,\zeta)
       \pm B(s,\zeta)\sqrt{1-\frac{\zeta}{\rho(s)}}.
\end{equation}
The normalization is $A(0,\zeta)=0$, $B(0,\zeta)=1$. The plus branch,
with the square root equal to one at $\zeta=0$, is the inverse germ
$v(s,0)=1$. Its only branch value within this disk is $\rho(s)$; that
branching is genuinely square-root. Its Taylor series about $\zeta=0$
has radius of convergence exactly $|\rho(s)|$.
\end{theorem}
\begin{proof}
For $j=1,2$ define
\begin{equation}\label{eq:contour-moments}
 S_j(s,\zeta)=\frac1{2\pi i}\int_{|v|=2}
            v^j\frac{\Phi_v(s,v)}{\Phi(s,v)-\zeta}\dd v.
\end{equation}
The denominator never vanishes on the contour in the parameter domain,
so these are holomorphic functions. By the argument principle they equal
$v_1^j+v_2^j$, with multiplicities, even though the individual roots need
not be single-valued. Set
\[
 A=\frac{S_1}{2},\qquad
 \Delta=\frac{S_2}{2}-A^2.
\]
For distinct roots $\Delta=(v_1-v_2)^2/4$, and this identity extends
through a double root. A zero of $\Delta$ occurs precisely when both
roots coalesce, hence precisely at $\zeta=\rho(s)$.

For completeness, simplicity of this discriminant zero follows directly
from a local change of variable. Write
$\Phi(s,v)-\rho(s)=(v-\vc(s))^2G(s,v)$, where
$G(s,\vc(s))=\Phi_{vv}(s,\vc(s))/2\ne0$.
The coordinate $\eta=(v-\vc)\sqrt{G(s,v)}$ has nonzero derivative
at the critical point and therefore an analytic local inverse.
The two roots are obtained by setting $\eta=\pm\sqrt{\zeta-\rho}$.
Squaring their half-difference shows
$\Delta=2(\zeta-\rho)/\Phi_{vv}(s,\vc)+\OO((\zeta-\rho)^2)$.
Thus the zero is simple. It follows that
\[
 C(s,\zeta):=\frac{\Delta(s,\zeta)}{1-\zeta/\rho(s)}
\]
extends holomorphically and never vanishes. The extension is joint in
$(s,\zeta)$: a holomorphic function vanishing on
$\zeta=\rho(s)$ is divisible by $\zeta-\rho(s)$, for example by integrating
its $\zeta$ derivative along the segment from $\rho(s)$ to $\zeta$ in a
local neighborhood. Nonvanishing at the zero follows from simplicity.

The parameter polydisk is simply connected, so the nonvanishing function
$C$ has a holomorphic square root $B$. At $s=0$ the roots are
$\pm\sqrt{1-2\zeta}$, hence $C(0,\zeta)=1$; choose $B=1$ there.
This proves \eqref{eq:atlas}. At $\zeta=0$ the plus branch agrees with
$v=1$ at $s=0$, and remains that simple root by holomorphic continuation.

On $|\zeta|<|\rho(s)|$ the square root normalized at zero is single-valued
and holomorphic, proving the lower bound on the Taylor radius. At
$\zeta=\rho(s)$ its coefficient $B(s,\rho(s))$ is nonzero, so this inverse
has a nonremovable square-root singularity there. The opposite inequality
follows. The contour argument also excludes any additional local root
collision or singularity of this two-root inverse inside $|\zeta|<3/4$.
\end{proof}

\begin{corollary}\label{cor:z-radius}
For $0<|s|<1/100$, the series \eqref{eq:sectors} at $w=-1+s$ has radius
of convergence exactly $|\zc(s)|$. This is a statement about the
independent variable $z$, with $w$ fixed.
\end{corollary}

\section{The moving discriminant and the nearby branch points}
\label{sec:critical}

\subsection{Critical-point expansions}
Equation \eqref{eq:k} is both an exact definition and a recursive method.
Because the linear coefficient of $k$ is $-1$, the coefficient of each
successive power of $s$ in $\tc$ is determined by the already known ones.
The result begins
\begin{align}
 \tc(s)={}&\frac{s^2}{2}+\frac{s^3}{3}-\frac{s^4}{8}
 -\frac{3s^5}{10}+\frac{s^6}{48}
 +\frac{323s^7}{840}+\frac{1123s^8}{5760}+\OO(s^9),
 \label{eq:tc-series}\\
 \rho(s)={}&\frac12+\frac s3+\frac{5s^2}{8}
 +\frac{7s^3}{10}+\frac{19s^4}{24}
 +\frac{361s^5}{420}+\OO(s^6).
 \label{eq:rho-series}
\end{align}
The second series follows by substituting the first in
\eqref{eq:rho-exact}. In particular,
\begin{equation}\label{eq:zc-series}
 \zc(s)=\frac{s^2}{2}+\frac{s^3}{3}+\frac{5s^4}{8}
                  +\frac{7s^5}{10}+\OO(s^6).
\end{equation}
The first omitted correction in the deepest-pole model is therefore
already cubic in $s$. Although smaller than the quadratic term as a
function value, it is not uniformly negligible after taking the square
root near its zero.

The logarithm of the normalized discriminant is also useful:
\begin{equation}\label{eq:log-rho}
 \log(2\rho(s))=\frac{2s}{3}+\frac{37s^2}{36}
                    +\frac{539s^3}{810}+\OO(s^4).
\end{equation}
All displayed coefficients in \eqref{eq:tc-series}--\eqref{eq:log-rho}
are checked independently by exact rational symbolic substitution in the
accompanying script.

\subsection{Branch points when the perturbation parameter is fixed}
The equation for branch values in the $(s,z)$ plane is
$z=\zc(s)=s^2\rho(s)$. It can be inverted without an ambiguous formal
choice of a square root. Define
\[
 U(s)=s\sqrt{2\rho(s)},\qquad U'(0)=1,
\]
using the holomorphic square root equal to one at $s=0$. Let $J$ be its
local holomorphic inverse.

\begin{proposition}\label{prop:branch-points}
For every sufficiently small nonzero $z$, the two nearby branch points
in the $w$ plane of the independent-parameter equation are
\begin{equation}\label{eq:branch-points}
 w_\pm(z)=-1+J(\pm\sqrt{2z}),
\end{equation}
where
\begin{equation}\label{eq:J}
 J(u)=u-\frac{u^2}{3}-\frac{25u^3}{72}
               +\frac{137u^4}{540}+\frac{6529u^5}{17280}+\OO(u^6).
\end{equation}
They are simple square-root branch points and coalesce at $w=-1$ as
$z\to0$. There are exactly two such points in a sufficiently small fixed
neighborhood of $w=-1$.
\end{proposition}
\begin{proof}
The equation $z=s^2\rho(s)$ is exactly $U(s)^2=2z$. Since $U$ is
biholomorphic near zero, it has precisely the two solutions in
\eqref{eq:branch-points}. Series substitution verifies \eqref{eq:J}.
For nonzero small $s$, $\zc'(s)=s+\OO(s^2)\ne0$, so the discriminant
zero is transverse to the fixed-$z$ slice. The factor $B$ in
Theorem~\ref{thm:atlas} is nonzero. Multiplication by $s$ in
$\delta=s(v-1)$ does not remove the square-root singularity because
$s\ne0$ at these points.
\end{proof}

This proposition locates actual local branch points of the two-variable
family, not merely zeros of a truncated quadratic surrogate. The
qualification in Section~\ref{sec:qualification} remains in force: it is
not a global classification for the map with coupling fixed at one.

\section{Polynomial profiles with an error bound at the fold}
\label{sec:profiles}

\subsection{The regular sheet coordinate}
Substitute $\zeta=\rho(s)(1-q^2)$ into the atlas and define
\begin{equation}\label{eq:V}
 V(s,q)=A\bigl(s,\rho(s)(1-q^2)\bigr)
       +qB\bigl(s,\rho(s)(1-q^2)\bigr).
\end{equation}
Changing $q$ to $-q$ gives the other sheet. Since
$|\rho(s)(1-q^2)|<0.504\cdot1.25<3/4$ for $|q|\le1/2$, this expression
is holomorphic for $|s|<1/100$, $|q|\le1/2$, and
\begin{equation}\label{eq:V-bound}
 |V(s,q)|<2,\qquad V(0,q)=q.
\end{equation}
At $q=0$ it equals the double root $\tc(s)/s$. The coalescence has become
an ordinary point of the $(s,q)$ parametrization.

\begin{theorem}[Polynomial profiles and uniform remainder]
\label{thm:profiles}
There are polynomials $V_j(q)\in\Q[q]$ such that
\begin{equation}\label{eq:profile-series}
 V(s,q)=\sum_{j\ge0}s^jV_j(q),\qquad
 V_0(q)=q,\qquad \deg V_j\le j+1.
\end{equation}
For $j\ge1$, $V_j(1)=0$ as a polynomial identity. The first five are
\begin{align}
 V_0(q)&=q,\nonumber\\
 V_1(q)&=\frac12+\frac q3-\frac{5q^2}{6},\nonumber\\
 V_2(q)&=\frac13-\frac{49q}{72}-\frac{5q^2}{9}
                                      +\frac{65q^3}{72},\label{eq:profiles}\\
 V_3(q)&=-\frac18-\frac{799q}{1080}+\frac{9q^2}{8}
                     +\frac{65q^3}{72}-\frac{157q^4}{135},\nonumber\\
 V_4(q)&=-\frac3{10}+\frac{22739q}{51840}+\frac{31q^2}{18}
          -\frac{3373q^3}{1728}-\frac{628q^4}{405}
          +\frac{28369q^5}{17280}.\nonumber
\end{align}
For $N\ge0$, put
\begin{equation}\label{eq:delta-N}
 \delta_N(s,q)=s\left(\sum_{j=0}^{N}s^jV_j(q)-1\right).
\end{equation}
Then for $|s|<1/100$ and $|q|\le1/2$,
\begin{equation}\label{eq:uniform-error}
 \left|s(V(s,q)-1)-\delta_N(s,q)\right|
 \le\frac{2|s|(100|s|)^{N+1}}{1-100|s|}.
\end{equation}
The estimate is valid at the branch point $q=0$ and is independent of the
sheet once its sign of $q$ has been chosen.
\end{theorem}
\begin{proof}
Holomorphy gives a convergent $s$ expansion with coefficients initially
holomorphic in $q$. We prove their stronger polynomial property from
\[
 \Phi(s,V(s,q))=\rho(s)(1-q^2).
\]
Write $\Phi(s,v)=\sum_{i\ge0}s^i\Phi_i(v)$. Equation
\eqref{eq:H-series} and the exponential factor show that
$\Phi_i\in\Q[v]$ and $\deg\Phi_i\le i+2$, with
$\Phi_0(v)=(1-v^2)/2$.

Suppose $V_0,\ldots,V_{j-1}$ are known. At order $s^j$ the only occurrence
of $V_j$ is $\Phi_0'(q)V_j=-qV_j$. Let $R_j(q)$ be the coefficient of
$s^j$ in
\begin{equation}\label{eq:profile-recursion}
 \Phi\left(s,\sum_{\ell=0}^{j-1}s^\ell V_\ell(q)\right)
                     -\rho(s)(1-q^2).
\end{equation}
Thus $qV_j=R_j$. By induction $R_j$ is a rational polynomial. For a term
of degree $d\le i+2$ in $\Phi_i$, the coefficient at total order $j$
after substitution has degree at most $d+(j-i)\le j+2$. The subtraction
on the right also obeys this bound. Hence $\deg R_j\le j+2$.

The analytically constructed coefficient $V_j$ is regular at $q=0$.
Consequently $R_j(0)=0$ and the polynomial $R_j$ is divisible by $q$.
It follows that $V_j\in\Q[q]$ and $\deg V_j\le j+1$. This argument is
important: formal division by $q$ alone would not justify regularity at
the fold. The exact atlas supplies that missing fact.

Near $q=1$ the same expression \eqref{eq:V} is defined for small $s$,
because its argument $\zeta$ is near zero. At $q=1$ it equals the root
$v=1$, proving $V_j(1)=0$. This is a statement about the coefficient
polynomials; the quantitative error domain asserted here remains
$|q|\le1/2$. Direct execution of the recurrence gives
\eqref{eq:profiles}.

Finally, \eqref{eq:V-bound} and Cauchy's estimate on circles $|s|=r'<r$
give $|V_j(q)|\le2(r')^{-j}$ uniformly for $|q|\le1/2$. Let $r'\uparrow r$,
sum the geometric tail, and multiply by $|s|$. This proves
\eqref{eq:uniform-error} for every $N$.
\end{proof}

\begin{remark}[What the error bound does and does not require]
The estimate concerns an exact input value of $q$, equivalently an exact
moving critical value $\zc(s)$. It does not silently include errors in
computing $\zc$, nor does it include a separate asymptotic truncation of
an exact logarithmic core at infinity. Errors in those quantities need
additional propagation bounds, especially when $q$ is small. A
residual-based bound valid at a fold is given in Section~\ref{sec:residual}.
\end{remark}

\section{Failure of the unshifted model and fractional actions}
\label{sec:failure}

\subsection{A sharp diagnostic path}
\begin{proposition}\label{prop:failure}
Let $s\to0^+$ and set $z=s^2/2$. On the sheet with positive $q$,
\begin{equation}\label{eq:failure}
 \delta=-s+\sqrt{\frac23}\,s^{3/2}+\frac{s^2}{2}
                        +\OO(s^{5/2}).
\end{equation}
The deepest-pole sum \eqref{eq:old-sqrt} equals $-s$ on this path.
Its error is therefore asymptotic to $\sqrt{2/3}\,s^{3/2}$ and is not
$\OO(s^2)$.
\end{proposition}
\begin{proof}
Equation \eqref{eq:rho-series} gives
\[
 q^2=1-\frac1{2\rho(s)}
       =\frac23s+\frac{29}{36}s^2+\OO(s^3),
 \qquad
 q=\sqrt{\frac23}s^{1/2}\bigl(1+\OO(s)\bigr).
\]
The polynomial profiles give
$\delta=s(q-1)+s^2(1/2+q/3-5q^2/6)+\OO(s^3)$ uniformly as $q\to0$.
Substitution proves \eqref{eq:failure}. The leading nonzero error divided
by $s^2$ diverges, proving the last assertion.
\end{proof}

This does not contradict the predecessor's coefficientwise pole theorem.
At each fixed $n$, all less singular Laurent terms disappear under the
normalization $s^{2n-1}$. On a path where $z/s^2$ approaches $1/2$, however,
infinitely many sectors are relevant and those discarded corrections move
the singularity itself. The interchange of these two operations is the
invalid step.

\subsection{An explicit convergent critical transseries}
Choose a real $a>0$ and specialize
\[
 s=e^{-aX},\qquad z=\frac12e^{-2aX},\qquad X\to+\infty.
\]
The plus branch has the convergent exponential transseries
\begin{equation}\label{eq:half-action}
 \delta(X)=-e^{-aX}+\sqrt{\frac23}\,e^{-3aX/2}
                       +\frac12e^{-2aX}+\OO(e^{-5aX/2}).
\end{equation}
Convergence, not just formal asymptotic validity, follows because $q$ is
$s^{1/2}$ times a holomorphic unit and $V$ is holomorphic in $(s,q)$.
After setting $r=s^{1/2}$, the whole expression is an ordinary convergent
power series in $r$ near zero.

On the exact critical curve $z=\zc(s)$, by contrast, $q=0$ and
\[
 \delta=\tc(s)-s=-s+\frac12s^2+\frac13s^3-\frac18s^4+\cdots.
\]
There are only integer multiples of the action $a$ on this curve. The
fractional action in \eqref{eq:half-action} is caused by detuning from the
moving critical curve, not by a failure of the analytic parametrization.

\subsection{How much discriminant accuracy is needed?}
Let $\zc^{[M]}(s)$ be the Taylor polynomial of $\zc(s)$ through degree
$M\ge2$, and suppose
\[
 \zc(s)-\zc^{[M]}(s)=c_{M+1}s^{M+1}+\OO(s^{M+2}),
 \qquad c_{M+1}\ne0.
\]
On the path $z=\zc^{[M]}(s)$ one has
\[
 q^2=2c_{M+1}s^{M-1}(1+\OO(s)).
\]
Since the half-difference between the two roots $\delta_\pm$ is
$sqB(s,\zeta)$ and $B\to1$, it follows that
\begin{equation}\label{eq:truncated-discriminant}
 \frac{\delta_+-\delta_-}{2}
 \sim\sqrt{2c_{M+1}}\,s^{(M+1)/2},
\end{equation}
with the chosen square-root convention. This statement concerns splitting
about the analytic center, so center shifts do not contaminate it.

For $M=2$, the coefficient is $c_3=1/3$ and the exponent is $3/2$.
For $M=3$, $c_4=5/8$ and the splitting is asymptotic to
$(\sqrt5/2)s^2$. This gives a simple accuracy rule: losing one order in
the critical value can lose half that order in the inverse near the fold.

\section{Uniform high-sector asymptotics and a crossover law}
\label{sec:crossover}

\subsection{The singular amplitude}
Let
\begin{equation}\label{eq:beta-gamma}
 \beta(s)=B(s,\rho(s))=\partial_qV(s,0),\qquad
 \gamma(s)=-\rho(s)B_\zeta(s,\rho(s))=\coeff{q}{3}V(s,q).
\end{equation}
The second identity for $\gamma$ follows by substituting
$\zeta=\rho(s)(1-q^2)$ into \eqref{eq:V}; the contribution from $A$ is
even in $q$.

\begin{lemma}\label{lem:beta}
The amplitude $\beta$ is holomorphic and nonzero for $|s|<1/100$,
with $\beta(0)=1$ and
\begin{equation}\label{eq:beta-exact}
 \beta(s)^2=
 \frac{2\tc(s)(1-\tc(s))}{s^2(1+\tc(s)-\tc(s)^2)}.
\end{equation}
In particular,
\begin{equation}\label{eq:beta-series}
 \beta(s)=1+\frac s3-\frac{49s^2}{72}-\frac{799s^3}{1080}
                         +\frac{22739s^4}{51840}+\OO(s^5),
 \qquad \gamma(s)=\OO(s^2).
\end{equation}
\end{lemma}
\begin{proof}
Expand $\Phi$ in $v$ at the critical point. If
$v=\vc+\beta q+\OO(q^2)$ and
$\zeta=\rho(1-q^2)$, comparison of quadratic terms gives
\[
 \beta^2=-\frac{2\rho}{\Phi_{vv}(s,\vc)}.
\]
Writing $E=e^{\tc-s}/(1-s)$, differentiation of \eqref{eq:Phi} and use of
\eqref{eq:k} yield
\[
 \Phi_{vv}(s,\vc)=E k'(\tc),\qquad
 k'(t)=-\frac{1+t-t^2}{(1-t)^2}.
\]
Combine this with \eqref{eq:rho-exact} to obtain
\eqref{eq:beta-exact}. The analytic choice of its square root is fixed
by $\beta(0)=1$. Equation \eqref{eq:beta-series} also follows by taking
the coefficients of $q$ and $q^3$ in \eqref{eq:profiles}.
\end{proof}

\subsection{A uniform coefficient lemma}
We record the elementary transfer argument needed below. It is a
specialized version of the coefficient-transfer principle in
\cite{flajolet-odlyzko}, with the uniformity made explicit.

\begin{lemma}\label{lem:transfer}
Suppose $\mathcal A_s(Z)$ and $\mathcal B_s(Z)$ are holomorphic in
$|Z|<R$ for some $R>1$, uniformly bounded on smaller closed disks, with
$s$ in a compact parameter set. Suppose
\[
 f_s(Z)=\mathcal A_s(Z)+\mathcal B_s(Z)(1-Z)^{1/2}
\]
as germs at zero. Set
$\beta_s=\mathcal B_s(1)$ and
$\gamma_s=-\mathcal B_s'(1)$. Then, uniformly in $s$,
\begin{equation}\label{eq:transfer}
 \coeff{Z}{n}f_s(Z)
 =-\frac{\beta_s}{2\sqrt\pi}n^{-3/2}
   +\left(-\frac{3\beta_s}{16\sqrt\pi}
          +\frac{3\gamma_s}{4\sqrt\pi}\right)n^{-5/2}
   +\OO(n^{-7/2}).
\end{equation}
An expansion to every algebraic order in $1/n$ follows from higher
Taylor coefficients of $\mathcal B_s$ at one.
\end{lemma}
\begin{proof}
The binomial coefficient identity, for the half-integers used here, is
\begin{equation}\label{eq:gamma-coeff}
 \coeff{Z}{n}(1-Z)^\alpha
 =\frac{\Gamma(n-\alpha)}{\Gamma(-\alpha)\Gamma(n+1)}.
\end{equation}
It follows either by multiplying the factors of the generalized binomial
coefficient or by induction in $n$. The gamma-ratio expansion gives
\begin{align*}
 \coeff{Z}{n}(1-Z)^{1/2}
   &=-\frac1{2\sqrt\pi}n^{-3/2}
                          \left(1+\frac3{8n}+\OO(n^{-2})\right),\\
 \coeff{Z}{n}(1-Z)^{3/2}
   &=\frac3{4\sqrt\pi}n^{-5/2}(1+\OO(n^{-1})).
\end{align*}
The needed constants are also obtained directly from the first two terms
of Stirling's formula; no parameter occurs in these scalar estimates.

Taylor division at $Z=1$ gives
\[
 \mathcal B_s(Z)=\beta_s+\gamma_s(1-Z)+(1-Z)^2\mathcal C_s(Z),
\]
where $\mathcal C_s$ is holomorphic on the same disk and uniformly bounded
on every smaller disk. To justify the coefficient remainder, fix
$1<R_1<R$. Its coefficients satisfy
$|\coeff{Z}{j}\mathcal C_s|\le M R_1^{-j}$ uniformly in $s$.
The coefficients $d_k$ of $(1-Z)^{5/2}$ satisfy
$|d_k|\le M'(k+1)^{-7/2}$, after enlarging $M'$ for finitely many small
indices. The convolution is bounded by
\[
 MM'\sum_{j=0}^nR_1^{-j}(n-j+1)^{-7/2}.
\]
For $j\le n/2$ this is $\OO(n^{-7/2})$ by a geometric sum. For
$j>n/2$ it is exponentially small, hence has the same bound. The analytic
term $\mathcal A_s$ likewise has exponentially small coefficients.
This proves \eqref{eq:transfer} with a uniform remainder. Repeating
Taylor division to any fixed order gives the final assertion by exactly
the same convolution estimate.
\end{proof}

\subsection{Application to the normalized sector functions}
\begin{theorem}[Uniform large-sector expansion]\label{thm:large-n}
Uniformly for $|s|\le1/200$, as $n\to\infty$,
\begin{equation}\label{eq:large-n}
 a_n(s)=-\frac{\beta(s)}{2\sqrt\pi}\rho(s)^{-n}n^{-3/2}
 \left[1+\frac{1}{n}
      \left(\frac38-\frac{3\gamma(s)}{2\beta(s)}\right)
       +\OO(n^{-2})\right].
\end{equation}
The expansion extends to every fixed algebraic order, with coefficients
holomorphic in $s$ and a remainder uniform on this closed disk.
\end{theorem}
\begin{proof}
Put $Z=\zeta/\rho(s)$. Theorem~\ref{thm:atlas} gives
\[
 v(s,\rho(s)Z)
 =\mathcal A_s(Z)+\mathcal B_s(Z)(1-Z)^{1/2},
\]
where $\mathcal A_s(Z)=A(s,\rho(s)Z)$ and
$\mathcal B_s(Z)=B(s,\rho(s)Z)$.
By Lemma~\ref{lem:rho-enclosure}, these functions are holomorphic in
$|Z|<7/5$ since $0.504\cdot7/5<3/4$. Restricting $|s|\le1/200$
leaves a compact subset of the $s$ domain, so the uniform boundedness
hypotheses of Lemma~\ref{lem:transfer} hold on smaller $Z$ disks.

The coefficient of $Z^n$ is $\rho(s)^n a_n(s)$ for $n\ge1$. Apply the
lemma and factor out the leading term. The division is uniform because
$\beta$ is nonzero and therefore bounded away from zero on the compact
$s$ disk. The all-order statement follows from the last assertion of the
lemma and holomorphic dependence of every local Taylor coefficient.
\end{proof}

The theorem includes $s=0$, where the coefficients are exactly
\begin{equation}\label{eq:a0-gamma}
 a_n(0)=-\frac{2^{n-1}\Gamma(n-1/2)}{\sqrt\pi\,\Gamma(n+1)}.
\end{equation}
For nonzero $s$, multiplication by $s^{1-2n}$ gives the corresponding
uniformly normalized asymptotic statement for the original rational
functions $b_n(-1+s)$.

\begin{theorem}[Critical/large-sector crossover]\label{thm:crossover}
For every $T>0$, uniformly for complex $|\tau|\le T$, one has
\begin{equation}\label{eq:crossover}
 \frac{a_n(\tau/n)}{a_n(0)}
 =e^{-2\tau/3}
 \left[1+\frac{\tau/3-37\tau^2/36}{n}
                       +\OO_T(n^{-2})\right]
 \qquad(n\to\infty).
\end{equation}
In particular, the limit is $e^{-2\tau/3}$, not one unless this
exponential happens to equal one.
\end{theorem}
\begin{proof}
For $n$ sufficiently large in terms of $T$, $s=\tau/n$ lies in the disk
of Theorem~\ref{thm:large-n}. Dividing that theorem by its value at $s=0$
and using \eqref{eq:beta-series} gives
\[
 \frac{a_n(s)}{a_n(0)}
 =\beta(s)(2\rho(s))^{-n}\bigl(1+\OO_T(n^{-2})\bigr).
\]
Indeed the coefficient of $1/n$ in \eqref{eq:large-n} is $3/8+\OO(s^2)$,
so its change contributes no larger error; the uniform remainders suffice.
The logarithm in \eqref{eq:log-rho} is the analytic one vanishing at zero.
Hence
\begin{align*}
 \beta(\tau/n)&=1+\frac{\tau}{3n}+\OO_T(n^{-2}),\\
 (2\rho(\tau/n))^{-n}
 &=\exp\left(-\frac{2\tau}{3}-\frac{37\tau^2}{36n}
                          +\OO_T(n^{-2})\right).
\end{align*}
Multiplication proves \eqref{eq:crossover}.
\end{proof}

\subsection{Why the exponential crossover is structurally stable}
The phenomenon is not tied to the particular coefficients above. Suppose
an analytic family has an atlas of the same form, with a unique dominant
simple square-root singularity at $\rho(s)$, nonzero amplitude $\beta(s)$,
and a common continuation disk after normalizing the singularity to one.
The proof of Theorem~\ref{thm:large-n} then gives
\[
 \frac{a_n(s)}{a_n(0)}\sim
 \frac{\beta(s)}{\beta(0)}
       \left(\frac{\rho(s)}{\rho(0)}\right)^{-n}
\]
uniformly for small $s$ as $n\to\infty$.
If
$\rho(s)/\rho(0)=1+\kappa s^k+\OO(s^{k+1})$, with
$\kappa\ne0$ and $k\ge1$, then any regime with $s\to0$ and
$ns^k\to\tau$ yields the limit $e^{-\kappa\tau}$.
This follows by taking an analytic logarithm; $ns^{k+1}\to0$ in such a
regime. It is a conditional universality statement with the analytic
continuation and nonvanishing hypotheses stated explicitly, not a claim
about arbitrary formal transseries.

\section{Analytic folds with finitely many exponential actions}
\label{sec:actions}

The previous example produces fractional exponential actions along a
critical path. We now isolate the general analytic mechanism and address
resonances without pretending that independent formal multi-indices
remain distinct after exponential specialization.

\subsection{Exact quadratic preparation without an unproved normal form}
\begin{lemma}[Local analytic fold preparation]\label{lem:fold}
Let $F(u,\bm q)$ be holomorphic near $(0,\bm0)\in\C\times\C^d$ and suppose
\begin{equation}\label{eq:fold-hyp}
 F(0,\bm0)=F_u(0,\bm0)=0,\qquad F_{uu}(0,\bm0)\ne0.
\end{equation}
After shrinking the neighborhood, there are holomorphic functions
$A(\bm q)$, $D(\bm q)$ and a nonvanishing holomorphic function
$\mathcal U(u,\bm q)$ such that
\begin{equation}\label{eq:fold-preparation}
 F(u,\bm q)=\mathcal U(u,\bm q)
             \bigl((u-A(\bm q))^2-D(\bm q)\bigr),
 \qquad A(\bm0)=D(\bm0)=0.
\end{equation}
The two nearby roots, with multiplicity, are
$u_\pm=A(\bm q)\pm\sqrt{D(\bm q)}$.
\end{lemma}
\begin{proof}
Choose a sufficiently small circle $|u|=r$ on which $F(u,\bm0)$ is
nonzero and encloses no zero other than its double zero at zero.
For sufficiently small $\bm q$, Rouch\'e's theorem preserves the count
of two zeros. Define the first two root moments by the integrals of
$u^jF_u/F$ over this circle, exactly as in
\eqref{eq:contour-moments}. Their symmetric combinations define
$A$ and $D$, and the monic quadratic
$P(u,\bm q)=(u-A)^2-D$ has exactly these two zeros with their
multiplicities.

For each fixed $\bm q$, $F/P$ is holomorphic within the circle after
removing its apparent zero-over-zero singularities. Joint holomorphy
can be verified without choosing the individual roots: for $|u|<r$ set
\[
 \mathcal U(u,\bm q)=\frac1{2\pi i}\int_{|\xi|=r}
  \frac{F(\xi,\bm q)}{P(\xi,\bm q)(\xi-u)}\dd\xi.
\]
The denominator on the contour is nonzero, so this is jointly holomorphic.
For fixed $\bm q$, Cauchy's formula identifies it with $F/P$ where the
quotient is already defined, and the equality extends everywhere inside.
At the origin its value is $F_{uu}(0,\bm0)/2\ne0$; shrink the neighborhood
once more to make it nonvanishing. This proves the factorization and
root formula.
\end{proof}

The lemma is the classical quadratic preparation mechanism, included
with a proof to make the later support construction independent of any
unstated normal-form theorem. No novelty is claimed for analytic
quadratic preparation itself.

\subsection{Grouping equal actions before taking a square root}
Fix real $\mu_1,\ldots,\mu_d>0$ and specialize
\begin{equation}\label{eq:action-specialization}
 q_i(X)=e^{-\mu_iX},\qquad X\to+\infty.
\end{equation}
Let
\begin{equation}\label{eq:S}
 S=\{\bm n\cdot\bm\mu:\bm n\in\N^d\}.
\end{equation}
This monoid is locally finite: for $\bm n\cdot\bm\mu\le M$, each
coordinate satisfies $n_i\le M/\mu_i$, so there are only finitely many
possible multi-indices. In particular, each action has only finitely
many representations by such a multi-index.

Expand the discriminant of Lemma~\ref{lem:fold} as
\[
 D(\bm q)=\sum_{\bm n\in\N^d}d_{\bm n}\bm q^{\bm n}.
\]
The expansion is absolutely convergent in a small polydisk. Under
\eqref{eq:action-specialization}, combine terms with the same action:
\begin{equation}\label{eq:grouped-D}
 D(X)=\sum_{\nu\in S}d_\nu e^{-\nu X},\qquad
 d_\nu=\sum_{\bm n\cdot\bm\mu=\nu}d_{\bm n}.
\end{equation}
These finite sums are the relevant coefficients. They are independent
of how the original monomials representing a given action are listed.
Because $D(\bm0)=0$, the coefficient $d_0$ is zero.

\begin{theorem}[Resonance-aware action splitting]\label{thm:actions}
Under \eqref{eq:fold-hyp} and \eqref{eq:action-specialization}, exactly one
of the following alternatives holds.

If all grouped coefficients $d_\nu$ vanish, then $D(X)=0$ for all
sufficiently large $X$, and the two roots coincide there.

Otherwise there is a least $\lambda>0$ with $d_\lambda\ne0$. Then the
roots have absolutely convergent exponential transseries
\begin{equation}\label{eq:action-root}
 u_\pm(X)=A(X)\pm\sqrt{d_\lambda}\,e^{-\lambda X/2}
       \sum_{k\ge0}\binom{1/2}{k}R(X)^k,
 \quad
 R(X)=\sum_{\nu>\lambda}\frac{d_\nu}{d_\lambda}
                              e^{-(\nu-\lambda)X}.
\end{equation}
For either consistent choice of $\sqrt{d_\lambda}$,
\begin{equation}\label{eq:splitting}
 u_+(X)-u_-(X)\sim2\sqrt{d_\lambda}\,e^{-\lambda X/2}.
\end{equation}
The expansions in \eqref{eq:action-root}, including expansion of each
power $R^k$ and grouping equal output actions, are absolutely convergent
for all sufficiently large $X$.
\end{theorem}
\begin{proof}
Absolute convergence of the original multivariable Taylor expansion for
large $X$ justifies grouping by action. Local finiteness of $S$ implies
that every nonempty subset of $S$ has a least member: choose one member,
then minimize over the finite set of smaller members. Thus the second
alternative has a well-defined first action $\lambda$.

Factor $d_\lambda e^{-\lambda X}$ from \eqref{eq:grouped-D}. All exponents
remaining in $R$ are positive. More is needed than pointwise smallness:
we establish smallness of the absolute coefficient sum. Choose $X_0$
large enough that the original expansion converges absolutely. If the
tail is nonempty, it has a least positive gap
$\eta=\min\{\nu-\lambda:\nu>\lambda,\ d_\nu\ne0\}$.
Then
\begin{align*}
 \sum_{\nu>\lambda}\left|\frac{d_\nu}{d_\lambda}\right|
                       e^{-(\nu-\lambda)X}
 &\le e^{-\eta(X-X_0)}
       \sum_{\nu>\lambda}\left|\frac{d_\nu}{d_\lambda}\right|
                       e^{-(\nu-\lambda)X_0}\\
 &\longrightarrow0.
\end{align*}
When the tail is empty the conclusion is immediate.
For large $X$ this absolute sum is less than one. Therefore the binomial
series in \eqref{eq:action-root} converges absolutely even after all its
products are expanded: it is majorized by
$\sum_k|\binom{1/2}{k}|\|R\|_X^k$, with $\|R\|_X<1$.
It is exactly a square root of $1+R(X)$ tending to one. Substitution in
Lemma~\ref{lem:fold} proves the root formula and leading splitting.
If every grouped coefficient vanishes, the absolutely convergent sum
for $D(X)$ is identically zero, proving the first alternative.
\end{proof}

\subsection{A constructive finite positive support grid}
There is a useful stronger conclusion than mere local finiteness of the
resulting support. Although the differences $\nu-\lambda$ need not
belong to the original positive monoid $S$, they lie in a finitely
generated positive extension that can be constructed explicitly.

\begin{theorem}[Finite-grid support after a fold]\label{thm:grid}
In the nonzero-discriminant alternative of Theorem~\ref{thm:actions},
put $\mu_{\max}=\max_i\mu_i$ and define
\[
 E_\lambda=\{\bm n\in\N^d:\bm n\cdot\bm\mu>\lambda\}.
\]
Let $M_\lambda$ be its componentwise minimal elements. Then $M_\lambda$
is finite and is contained in the explicit box
\begin{equation}\label{eq:minimal-box}
 0\le m_i\le
 \left\lfloor\frac{\lambda+\mu_{\max}}{\mu_i}\right\rfloor
 \qquad(1\le i\le d).
\end{equation}
Define the positive numbers
$\beta_{\bm m}=\bm m\cdot\bm\mu-\lambda$ for
$\bm m\in M_\lambda$, and let $H$ be the additive monoid generated by
\begin{equation}\label{eq:grid-generators}
 \mu_1,\ldots,\mu_d,
 \quad \{\beta_{\bm m}:\bm m\in M_\lambda\}.
\end{equation}
Then
\begin{equation}\label{eq:grid-support}
 \supp A(X)\subseteq H,\qquad
 \supp\sqrt{D(X)}\subseteq\lambda/2+H,\qquad
 \supp u_\pm\subseteq H\cup(\lambda/2+H).
\end{equation}
In particular, adjoining $\lambda/2$ to the finite list
\eqref{eq:grid-generators} gives a single finite positive generating set
containing both root supports. If all $\mu_i$ lie in
$(a/N)\N$ for some $a>0$ and positive integer $N$, both supports lie in
$(a/(2N))\N$.
\end{theorem}
\begin{proof}
The set $E_\lambda$ is upward closed under the componentwise order.
If $\bm m$ is minimal and $m_i>0$, then
$(\bm m-\bm e_i)\cdot\bm\mu\le\lambda$, for otherwise minimality would
fail. Hence
$\bm m\cdot\bm\mu\le\lambda+\mu_i\le\lambda+\mu_{\max}$.
Every coordinate obeys \eqref{eq:minimal-box}, so the set of minimal
indices is finite. This also gives a direct finite search for it, once
$\lambda$ and exact comparisons among the actions are available.

For each $\bm n\in E_\lambda$, decrease coordinates while remaining in
$E_\lambda$ until no further decrease is possible. The process terminates
because the sum of coordinates decreases. It produces
$\bm m\in M_\lambda$ with $\bm m\le\bm n$. Therefore
\[
 \bm n\cdot\bm\mu-\lambda
 =\beta_{\bm m}+(\bm n-\bm m)\cdot\bm\mu\in H.
\]
Every positive gap appearing in $R$ has this form. Thus $\supp R\subseteq H$,
and every expanded power $R^k$ has support in $H$. The absolute
convergence proved in Theorem~\ref{thm:actions} justifies collecting
these powers; local finiteness also follows because $H$ has finitely many
strictly positive generators. Multiplication by $e^{-\lambda X/2}$ gives
the second inclusion in \eqref{eq:grid-support}.
The Taylor expansion of $A$ has support in $S\subseteq H$, proving the
other inclusions. If the original actions lie in $(a/N)\N$, then so do
$\lambda$ and all the positive differences $\beta_{\bm m}$; halving
$\lambda$ only requires the stated doubled denominator.
\end{proof}

The finite box is a sufficient construction, not an assertion that its
resulting generating list is minimal. A smaller list may be obtained
from the actual support of the discriminant. Also, identifying the first
nonzero grouped coefficient is not claimed to be an automatically
decidable operation for an arbitrary analytic input. When the nonzero
case is known and exact coefficient and action comparisons are available,
a search in increasing action terminates; the identically zero case may
require a separate identity proof.

\subsection{A cancellation that changes the first fractional action}
\begin{example}\label{ex:resonance}
Consider
\[
 F(u,q_1,q_2)=u^2-(q_1^2-q_2+q_1^3),
 \qquad (\mu_1,\mu_2)=(1,2).
\]
Before specialization, the terms $q_1^2$ and $-q_2$ are distinct monomials.
After specialization they have the same action and cancel exactly:
\[
 D(X)=e^{-2X}-e^{-2X}+e^{-3X}=e^{-3X}.
\]
Consequently $u_\pm=\pm e^{-3X/2}$. A rule that selected one monomial of
smallest action before combining resonances would incorrectly predict
splitting at action $1$ instead of $3/2$.
If the term $q_1^3$ is removed, the discriminant vanishes identically on
the specialized curve and the two roots coincide there, even though the
multivariable discriminant $q_1^2-q_2$ is not the zero function.
\end{example}

This example distinguishes an analytic identity in all perturbation
variables from an identity after restriction to an exponential curve.
The action theorem needs only the latter, and its grouping prescription
makes that distinction explicit.

\subsection{Relation to the broader transseries questions}
Theorems~\ref{thm:actions} and \ref{thm:grid} concern analytic germs in a
\emph{finite} collection of exponentially small parameters. Their
convergence follows from ordinary Taylor convergence and a binomial
majorant. They do not assert convergence for arbitrary Hahn series,
realization of unrestricted formal supports, Borel summability of
Gevrey series, or preservation of resurgent Stokes data. Those are
separate problems. The result is nevertheless more than a formal
square-root manipulation: it proves representation-independent handling
of resonances, absolute summability after all regroupings, and a finite
positive support grid with a concrete search bound.

\section{Residual certificates that remain valid at the fold}
\label{sec:residual}

A derivative-based inverse error bound usually divides by a quantity
that tends to zero at a critical point. Quadratic preparation gives a
replacement that is uniform across the collision.

\begin{proposition}[Root-set residual bound]\label{prop:residual}
Assume a factorization \eqref{eq:fold-preparation} is known on a domain
containing a candidate $\widetilde u$ and both local roots. Suppose
$|\mathcal U(\widetilde u,\bm q)|\ge m>0$, and put
\[
 r=F(\widetilde u,\bm q),\qquad R=|r|/m,
 \qquad d=\min\{|\widetilde u-u_+|,|\widetilde u-u_-|\}.
\]
Then
\begin{equation}\label{eq:residual}
 d\le\sqrt R.
\end{equation}
If $D(\bm q)\ne0$, the stronger combined estimate is
\begin{equation}\label{eq:residual-combined}
 d\le\min\left\{\sqrt R,\frac{R}{|\sqrt{D(\bm q)}|}\right\}.
\end{equation}
The square-root exponent in \eqref{eq:residual} is optimal at a double
root in general.
\end{proposition}
\begin{proof}
The preparation gives
\[
 |(\widetilde u-u_+)(\widetilde u-u_-)|\le R.
\]
The product is at least $d^2$, proving \eqref{eq:residual}. If the roots
are distinct, their separation is $2|\sqrt D|$. By the triangle
inequality at least one of the two candidate-to-root distances is at
least $|\sqrt D|$. Multiplying the smaller distance by this lower bound
proves $d\le R/|\sqrt D|$.
For $F(u)=u^2$, the exact distance to the root is $\sqrt{|F(u)|}$, so no
uniform replacement by a higher power of the residual is possible.
\end{proof}

The estimate certifies proximity to the \emph{root set}, not a preferred
sheet. A sheet-specific certificate requires continuation from a known
base point, an isolating region, or other branch information. This
qualification is necessary near a collision: the residual alone cannot
distinguish two very close roots.

A practical verification pipeline can therefore combine a polynomial
profile to obtain a candidate with a rigorous enclosure of the residual
and the unit factor. Away from the fold, the second term in
\eqref{eq:residual-combined} recovers linear error transport. At the fold,
the first term remains finite. The accompanying script does not implement
interval enclosures of the unit factor; the proposition gives the
mathematical certificate that such an implementation should check.

\section{Computation, reproducibility, and proof status}
\label{sec:computation}

\subsection{What the symbolic checks verify}
The package contains \texttt{verify.py}, using SymPy and mpmath. The
exact part checks the critical-point equation through degree eight,
$\rho$ through degree five, the first five profile polynomials, the
amplitude formula through degree four, the logarithm
\eqref{eq:log-rho}, and the branch-point reversion \eqref{eq:J}.
The rational inequalities used for the explicit domain are checked as
exact rational comparisons, not by a floating-point threshold test.

These tests can detect incorrect coefficients or arithmetic. They do not
prove the existence of the holomorphic atlas, all-order polynomiality,
the uniform coefficient theorem, or the finite-grid result. Those claims
are established by the arguments in the preceding sections. The script
also independently evaluates the first three known sector formulas
\eqref{eq:first-bn} as a consistency check on the normalization.

\subsection{Computing high sector indices without symbolic expression swell}
A scalar Lagrange formula provides an $\OO(n^2)$ arithmetic algorithm for
$a_n(s)$ at a fixed numerical value of $s$. Put $u=v-1$. Direct expansion
of \eqref{eq:Phi} gives
\begin{equation}\label{eq:J-unit}
 \Phi(s,1+u)=-\frac{e^{su}u}{(1-s)^2}J_s(u),\qquad
 J_s(u)=1+\sum_{j\ge1}
     \frac{(s/(1-s))^{j-1}}{(1-s)(j+1)}u^j.
\end{equation}
For $s=0$, the removable interpretation is $J_0(u)=1+u/2$.
To verify the formula directly, use
\[
 h(s(1+u))-h(s)
 =su+\log\left(1-\frac{s}{1-s}u\right)
\]
and separate its linear term in $u$.

Since $\Phi_v(s,1)\ne0$, ordinary scalar Lagrange inversion gives
\begin{equation}\label{eq:lagrange-numeric}
 a_n(s)=\frac{(-1)^n(1-s)^{2n}}n
         \coeff{u}{n-1}\left(e^{-nsu}J_s(u)^{-n}\right).
\end{equation}
This formula is also a purely formal identity in $u$ over the rational
function field in $s$; it does not require a convergence assumption to
extract a fixed coefficient.

Write $J_s(u)=\sum_{j\ge0}J_j u^j$, $J_0=1$, and
$J_s(u)^{-n}=\sum_{k\ge0}C_k u^k$. Differentiating
$C=J_s^{-n}$ gives $JC'=-nJ'C$. Comparing coefficients yields
\begin{equation}\label{eq:unit-recurrence}
 C_0=1,\qquad
 C_k=\frac1k\sum_{j=1}^k\bigl((1-n)j-k\bigr)J_jC_{k-j}.
\end{equation}
Computing $C_0,\ldots,C_{n-1}$ takes $\OO(n^2)$ arithmetic operations;
a final convolution with $e^{-nsu}$ evaluates
\eqref{eq:lagrange-numeric}. This is an operation count, not a claim of
quadratic bit complexity or numerical stability at arbitrary precision.
The supplied numerical calculations use 100 decimal digits.

\subsection{Numerical checks at and close to the branch point}
Table~\ref{tab:profiles} compares the order-$N=4$ profile truncation with
a high-precision root. At $q=0$ the reference root is the independently
computed $\tc(s)-s$. For $q\ne0$, the local root is computed directly
from \eqref{eq:Z}, initialized on the sheet selected by the profile.

\begin{table}[htbp]
\centering
\sisetup{exponent-mode=scientific,round-mode=none}
\begin{tabular}{@{}rrrr@{}}
\toprule
$s$ & $q$ & {Observed absolute error} & {Bound \eqref{eq:uniform-error}}\\
\midrule
\tableinput{data/profile_table.tex}
\bottomrule
\end{tabular}
\caption{Order-four profile checks, including the exact branch point and
an extremely close point. The bound column is rounded upward at the
shown precision. The observations are floating-point diagnostics, not
interval-arithmetic certificates.}
\label{tab:profiles}
\end{table}

The proved bound is intentionally conservative: its coefficient majorant
comes from a common complex domain and the simple root bound $|v|<2$,
not from optimized estimates for these particular real inputs. Its
purpose is to prove uniform validity through coalescence, rather than to
predict the actual error sharply.

\subsection{Numerical checks of the joint limit}
Table~\ref{tab:crossover} evaluates
$a_n(\tau/n)/a_n(0)$ by \eqref{eq:lagrange-numeric} and compares it with
both the limiting exponential and the first corrected expression in
Theorem~\ref{thm:crossover}. The decreasing corrected errors are
consistent with the proved $\OO_T(n^{-2})$ remainder.

\begin{table}[htbp]
\centering
\small
\begin{tabular}{@{}rrrrr@{}}
\toprule
$\tau$ & $n$ & $a_n(\tau/n)/a_n(0)$ & $e^{-2\tau/3}$
 & {Corrected error}\\
\midrule
\tableinput{data/crossover_table.tex}
\bottomrule
\end{tabular}
\caption{Large-sector crossover diagnostics. ``Corrected error'' is the
absolute difference from the full displayed right-hand side of
\eqref{eq:crossover} with its $\OO_T(n^{-2})$ term omitted.}
\label{tab:crossover}
\end{table}

The JSON output records the decimal precision, package versions, exact
expressions, and all test rows. The code writes its output relative to
its own directory, so running it from another working directory does not
silently overwrite unrelated data files.

\subsection{Formalization boundaries and a feasible first tranche}
No new Lean file is included or claimed verified. The repository's
existing Catalan module supplies useful finite algebra, but it does not
prove the complex analytic statements of this manuscript. The following
separation would keep a future formalization honest and manageable.

\begin{center}
\small
\begin{tabular}{@{}>{\raggedright\arraybackslash}p{0.26\textwidth}>{\raggedright\arraybackslash}p{0.33\textwidth}>{\raggedright\arraybackslash}p{0.32\textwidth}@{}}
\toprule
Layer & Mathematical content & Present status\\
\midrule
Finite algebra & Displayed rational coefficients, polynomial recurrences,
Catalan normalization & Exact symbolic checks; related Catalan core
already has a repository Lean module.\\[4pt]
Local analysis & Root count, holomorphic moments, moving discriminant,
uniform profile remainder & Proofs here; no new machine-checked proof.\\[4pt]
Coefficient asymptotics & Common continuation disk, half-integer transfer,
joint scaling limit & Proofs here and numerical diagnostics; not formalized.\\[4pt]
Action support & Resonance grouping, finite positive grid, absolute
summability & Proofs here; only the toy resonance is numerically checked.\\
\bottomrule
\end{tabular}
\end{center}

A useful first formalization target is the finite polynomial recurrence,
together with a theorem that an independently supplied analytic solution
forces the divisibility by $q$. A second target is the purely ordered
finite-box argument in Theorem~\ref{thm:grid}. These avoid claiming that
formal coefficient identities already imply analytic convergence. The
contour atlas and uniform transfer theorem are substantial later targets.

\section{Further research questions and concrete next steps}
\label{sec:questions}

The following questions distinguish genuine continuations from results
already proved above. They range from implementable certification tasks
to broader analytic and transseries problems.

\begin{question}[Sharp domains and practical certified evaluation]
How far can the explicit atlas domain be enlarged while retaining a
single simple critical value and a common contour enclosing precisely two
roots? Replace the conservative constants by interval-certified bounds,
and implement evaluation of $\tc$, $\rho$, the profiles, and the residual
certificate without losing reliability when $q$ is close to zero.
A useful deliverable is a branch-aware evaluator with an absolute error
guarantee, not merely a higher-order symbolic expansion.
\end{question}

\begin{question}[The fixed-coupling global continuation problem]
Determine the singularity geometry of the inverse of
$y+\log y+e^{-y}$ under its actual fixed coupling. Which critical points
and asymptotic values obstruct continuation of the positive real inverse?
How does the small-coupling chart here continue as $\lambda$ approaches
one? The relation $\lambda=we^wz$ must be enforced rather than treating
$z$ as independent at the end of a local argument.
\end{question}

\begin{question}[Two-stage truncation at infinity]
For the real positive inverse, replace the exact logarithmic core $w$ by
a truncated power--logarithmic expansion while also truncating the
exponential sectors. Prove a bound uniform in both truncation orders and
specify how they may grow with the large argument. The present theorem
is uniform near a critical auxiliary chart and does not settle the
predecessor's distinct two-stage uniformity question at infinity.
\end{question}

\begin{question}[Higher critical multiplicity]
For analytic equations with $F_u=\cdots=F_{u^{m-1}}=0$ and
$F_{u^m}\ne0$, construct quantitative $m$-sheet profiles analogous to
Theorem~\ref{thm:profiles}. Identify precisely which coefficients remain
polynomial after a moving-discriminant normalization and when genuinely
multivariate Newton polygons are unavoidable. Determine sharp support
grid enlargements beyond the simple denominator doubling of a fold.
\end{question}

\begin{question}[Competing singularities and new coefficient windows]
Extend the uniform large-sector theorem to a family in which two or more
critical values have comparable modulus, or themselves collide. Derive
joint limits that include interference among their coefficient
contributions. The common single-singularity continuation disk used here
will generally fail, so simply adding two fixed-parameter asymptotic
formulas is not sufficient.
\end{question}

\begin{question}[Resurgence beyond analytic perturbation parameters]
Replace convergent Taylor coefficients in the perturbation variables by
Gevrey or resurgent coefficient series. Under what explicit hypotheses
does moving-core inversion preserve Borel summability, and how are
singular locations and Stokes constants transformed? The absolute
convergence proof in Section~\ref{sec:actions} supplies none of these
conclusions for a divergent input and should not be used as a substitute
for them.
\end{question}

\begin{question}[Effective cancellation and minimal action grids]
For an effectively represented analytic discriminant, design certificates
for the first nonzero grouped action and for exact vanishing on a given
exponential curve. In the rationally dependent case, identify useful
algebraic classes where the resonance identities reduce to a finite
ideal-membership problem. In the incommensurate case, specify exact
comparison assumptions. Minimize the sufficient grid
\eqref{eq:grid-generators} without erasing cancellations or introducing
negative generators.
\end{question}

\begin{question}[Actions depending on the large parameter]
What remains true when $\mu_i=\mu_i(X)$ or when an external parameter
makes distinct actions approach one another? The positive gap in the
proof of Theorem~\ref{thm:actions} may cease to be uniform. Develop a
joint scaling theory that distinguishes a genuine resonance from a
near-resonance and identifies when a fixed finite support grid must be
replaced by a parameter-dependent description.
\end{question}

\begin{question}[Proof-producing critical inversion]
Build an implementation that emits a finite certificate for the root
count, critical-value enclosure, polynomial-profile residual, and
root-set error bound. A small trusted checker should reject unsupported
sheet choices and unproved nonvanishing assumptions. Begin with the
rational majorants and finite support-grid construction; then connect
these certificates to a Lean formalization of the analytic hypotheses.
\end{question}

\begin{question}[Gamma and Barnes cores]
Apply the method to analytic families genuinely underlying the
repository's Gamma and Barnes deepest-pole calculations. Determine the
moving critical values and corresponding joint large-sector limits.
Matching a quadratic coefficient recurrence is not enough: the analytic
family, its continuation domain, and the nonvanishing singular amplitude
must each be supplied and verified for those special functions.
\end{question}

\section{Conclusion}
The coefficientwise pole hierarchy and its formal square-root sum do
not by themselves provide a uniform inverse near a critical point.
For the mixed logarithmic--exponential sector equation, the missing
object is now explicit: a moving discriminant, an exact two-sheet
holomorphic atlas, and a convergent polynomial profile expansion with an
error bound valid at the collision itself. This resolves the predecessor's
stated local critical-transition question within its independent-parameter
family.

The same atlas controls a different limit that the fixed-sector pole
calculation does not see. High sector index amplifies the cubic drift of
the critical value, producing the nontrivial crossover
$e^{-2\tau/3}$ when $s=\tau/n$. The proof identifies both the leading
limit and its first finite-$n$ correction.

For analytic folds with several exponential inputs, the decisive datum
is the first nonzero \emph{grouped} discriminant action. Its half controls
branch splitting, while a constructive finite positive grid controls all
higher terms. Resonances can delay the first splitting action or remove
it altogether. The resulting series are absolutely convergent for the
analytic finite-action class considered here.

These conclusions do not solve general transseries realization,
unrestricted resurgence, or the global singularity geometry of the
fixed-coupling original map. They do provide a proved, quantitative
bridge from formal deepest poles to local analytic inversion, and a
concrete next set of questions with their remaining hypotheses visible.

\appendix
\section{Dependency map and source audit}
\label{app:audit}

The principal proof dependencies are
\[
 \text{explicit bounds}\ \Longrightarrow\ \text{two-root contour atlas}
 \ \Longrightarrow\
 \begin{cases}
   \text{uniform polynomial profiles},\\
   \text{uniform high-sector transfer},
 \end{cases}
\]
followed by the explicit critical-value expansion for the crossover.
The general action theorem uses a separate local fold-preparation lemma,
absolute Taylor convergence, and the finite positive-gap construction.
No general theorem of formal transseries composition or inversion is
assumed in either chain.

The repository was inspected at commit
\begin{center}
\small\texttt{20301601562cf84173f2b503b0dc9c3fbf21cd0c}.
\end{center}
The package README and \texttt{QuadraticCoreCatalan.lean} were read at that
pinned revision. The large canonical transseries source could not be
retrieved in full through the available repository fetch operation;
therefore this manuscript does not claim a complete audit of that
volume or the whole repository. The directly inspected status materials
are sufficient for the limited formalization distinctions made here.

The predecessor source \texttt{reversion\_and\_one\_exponential.tex} was
retrieved independently from the user's saved document library. Its
mixed-inverse equation, exact sector formula, leading-pole statement,
uniform real-positive sector bound, and critical-transition research
question were read directly. It is cited as a supplied research
manuscript, not as an externally peer-reviewed result. The main local
critical analysis in this article is proved from equation
\eqref{eq:delta}, so its validity does not require accepting any broader
claims of that predecessor.

\section{Reproduction instructions}
\label{app:reproduction}
The package is self-contained for compilation: all tables used by the
article are included. From the package directory, run
\begin{verbatim}
python -m pip install -r requirements.txt
python verify.py
latexmk -pdf -interaction=nonstopmode -halt-on-error article.tex
\end{verbatim}
The verification script has no network dependency after installation.
Its default run includes sector indices up to 400; \texttt{--quick}
omits those final large-index diagnostics. The symbolic checks are
identical in both modes. Re-running the script intentionally replaces
the package's own \texttt{data/verification.json} and table fragments;
it does not modify any repository or Library source.

The supplied PDF was compiled from this source and visually checked.
The package records a mathematical proof manuscript and reproducible
checks, not a claim of independent peer review, a formal verification
certificate, or established worldwide priority.

\begingroup
\small
\begin{thebibliography}{9}

\bibitem{repo-readme}
Vladimir Reshetnikov, \emph{ProveIt: Transseries\_And\_Inversion package
README}, inspected at commit
\texttt{20301601562cf84173f2b503b0dc9c3fbf21cd0c}, 29 September 2026.
\href{https://github.com/VladimirReshetnikov/ProveIt/blob/20301601562cf84173f2b503b0dc9c3fbf21cd0c/Analysis/FabiusFunction/docs/semi-formalized-research-frontiers/drafts/series-and-transseries/Transseries_And_Inversion/README.md}{Pinned repository source}.
The cited status concerns the distinctions between finite algebra,
assembled formal-series theorems, and analytic identification.

\bibitem{repo-catalan}
Vladimir Reshetnikov, ProveIt module
\emph{QuadraticCoreCatalan.lean}, same pinned revision.
\href{https://github.com/VladimirReshetnikov/ProveIt/blob/20301601562cf84173f2b503b0dc9c3fbf21cd0c/Analysis/FabiusFunction/Lean/FabiusFunction/QuadraticCoreCatalan.lean}{Pinned Lean source}.
Relevant volume labels are
\repoLabel{p6:prop:quadratic-core-catalan},
\repoLabel{p6:lem:quadratic-core}, and
\repoLabel{p6:thm:deepest-pole}; the module states its formalization boundaries.

\bibitem{predecessor}
\emph{Support-Controlled Reversion and the Exact One-Exponential
Substitution Group}, supplied research manuscript, 29 September 2026.
Source file \texttt{reversion\_and\_one\_exponential.tex}; retrieved from
the user's saved documents. In particular, the mixed inverse example,
its deepest-pole resummation, and Research question 11.8.
This is project-local predecessor material, not a claim of peer-reviewed
publication.

\bibitem{corless}
R.~M. Corless, G.~H. Gonnet, D.~E.~G. Hare, D.~J. Jeffrey, and D.~E. Knuth,
\emph{On the Lambert W function}, Advances in Computational Mathematics
\textbf{5} (1996), 329--359.
\href{https://doi.org/10.1007/BF02124750}{doi:10.1007/BF02124750}.
The branch-point expansion is discussed on pp.~350--351.

\bibitem{flajolet-odlyzko}
P. Flajolet and A. Odlyzko,
\emph{Singularity analysis of generating functions},
SIAM Journal on Discrete Mathematics \textbf{3}, no.~2 (1990), 216--240.
\href{https://algo.inria.fr/flajolet/Publications/FlOd90b.pdf}{Author-hosted article}.
The half-integer coefficient identity and its asymptotic expansion are
special cases of the coefficient formulas on p.~219.

\bibitem{edgar}
G.~A. Edgar, \emph{Transseries for beginners},
Real Analysis Exchange \textbf{35} (2010), 253--310.
\href{https://arxiv.org/abs/0801.4877v5}{arXiv:0801.4877v5}, revised 15 July 2009.

\end{thebibliography}
\endgroup
\end{document}
